% ============================================================================
% CHEATSHEET TEMPLATE
% ============================================================================
% Based on the ANA-G cheatsheet style
%
% Features:
% - 4-column landscape layout on A4
% - Colored boxes for main concepts and sub-concepts
% - Compact formatting optimized for information density
% - Mathematical typesetting support
%
% Usage:
% 1. Replace metadata (title, author, date)
% 2. Adjust number of columns if needed (see \Ncolumn)
% 3. Add your content using mainbox and subbox environments
% 4. Compile with pdflatex (run twice for proper layout)
% ============================================================================

\documentclass[a4paper, fontsize=7pt]{scrartcl}

% ============================================================================
% ENCODING & LANGUAGE
% ============================================================================
\usepackage[utf8]{inputenc}
% \usepackage[german]{babel}    % Uncomment for German hyphenation
% \usepackage[english]{babel}   % Uncomment for English hyphenation

% ============================================================================
% PAGE LAYOUT
% ============================================================================
% 4-column landscape layout with minimal margins
\usepackage[landscape, left=0.75cm, top=1cm, right=0.75cm, bottom=1.5cm, footskip=5pt]{geometry}
\usepackage{flowfram}
\ffvadjustfalse
\setlength{\columnsep}{1cm}

% Number of columns (adjust as needed: 3-4 recommended for A4)
\Ncolumn{4}

% Remove paragraph indentation
\setlength{\parindent}{0cm}

% ============================================================================
% COLORS & BOXES
% ============================================================================
\usepackage[most]{tcolorbox}

% Base style for all boxes
\tcbset{
  base/.style={
    boxrule=0.3mm,
    leftrule=1mm,
    left=1.75mm,
    arc=0mm,
    fonttitle=\bfseries,
    colbacktitle=black!10!white,
    coltitle=black,
    toptitle=0.75mm,
    bottomtitle=0.25mm,
    title={#1}
  }
}

% Main concept box (green accent)
% Use for: definitions, theorems, important formulas
\definecolor{maincolor}{rgb}{0.004, 0.5, 0.35}  % Nice green
\newtcolorbox{mainbox}[1]{
  colframe=maincolor,
  base={#1}
}

% Sub concept box (gray accent)
% Use for: examples, remarks, minor concepts
\newtcolorbox{subbox}[1]{
  colframe=black!20!white,
  base={#1}
}

% Alternative colors (uncomment to use)
% \definecolor{maincolor}{rgb}{0.2, 0.4, 0.8}   % Blue
% \definecolor{maincolor}{rgb}{0.8, 0.2, 0.2}   % Red
% \definecolor{maincolor}{rgb}{0.5, 0.2, 0.6}   % Purple

% ============================================================================
% MATHEMATICS
% ============================================================================
\usepackage{amsmath, amsthm, mathtools, amssymb}
\usepackage{amsfonts, amstext}
\allowdisplaybreaks  % Allow page breaks in equations

% Common math shortcuts (customize as needed)
\def\limn{\lim_{n\to \infty}}
\def\limxo{\lim_{x\to 0}}
\def\limxi{\lim_{x\to\infty}}
\def\limxn{\lim_{x\to-\infty}}
\def\sumk{\sum_{k=1}^\infty}
\def\sumn{\sum_{n=0}^\infty}

% Number sets
\def\N{\mathbb{N}}  % Natural numbers
\def\Z{\mathbb{Z}}  % Integers
\def\Q{\mathbb{Q}}  % Rationals
\def\R{\mathbb{R}}  % Reals
\def\C{\mathbb{C}}  % Complex

% Calculus
\def\dx{\,dx}
\def\dt{\,dt}
\def\dy{\,dy}
\newcommand{\ddx}[1]{\frac{d#1}{dx}}
\newcommand{\ddn}[2]{\frac{d#1}{d#2}}
\newcommand{\pdn}[2]{\frac{\partial #1}{\partial #2}}

% Big O notation
\def\oof{\mathcal{O}}

% ============================================================================
% TABLES
% ============================================================================
\usepackage{tabularx, multirow, booktabs}
\renewcommand*{\arraystretch}{1.5}  % More spacing in tables

% ============================================================================
% LISTS
% ============================================================================
\usepackage{enumitem}
\setitemize{itemsep=0.5pt, topsep=2pt}
\setenumerate{itemsep=0.75pt, topsep=2pt}

% ============================================================================
% GRAPHICS & TIKZ
% ============================================================================
\usepackage{graphicx}
\usepackage{tikz}
\usetikzlibrary{
  positioning,
  calc,
  fit,
  matrix,
  decorations,
  angles,
  quotes,
  babel,
  3d,
  arrows.meta,
  shadings,
  shapes.geometric
}

\tikzset{
  >=stealth,
  % Add custom TikZ styles here
}

% ============================================================================
% HYPERLINKS
% ============================================================================
\usepackage{hyperref}
\hypersetup{
  colorlinks=true,
  linkcolor=maincolor,
  urlcolor=maincolor
}

% ============================================================================
% ADDITIONAL PACKAGES
% ============================================================================
\usepackage{centernot}   % For \centernot command
\usepackage{marvosym}    % Additional symbols
\usepackage{wasysym}     % Additional symbols

% ============================================================================
% DOCUMENT METADATA
% ============================================================================
\title{DMATH-ALGO Cheatsheet}
\author{Discrete Mathematics \& Algorithms}
\date{Version 2024}

% Table of contents depth
\setcounter{tocdepth}{4}
\setcounter{secnumdepth}{4}

% ============================================================================
% DOCUMENT START
% ============================================================================
\begin{document}

\maketitle

% ============================================================================
% SECTION 1: BOOLEAN
% ============================================================================
\section{Boolean}

\begin{mainbox}{Definition: Boolean}
Ein boolescher Wert ist entweder \texttt{True} oder \texttt{False}.

\textbf{Darstellung:}
\begin{itemize}
    \item 1 fur \texttt{True}, 0 f ¨ ur \texttt{False} ¨
    \item Bit: grundlegendste Einheit von Daten
\end{itemize}
\end{mainbox}

\begin{mainbox}{Boolesche Operatoren}
\textbf{Einstelliger Operator not:}
\begin{center}
\begin{tabular}{c|c}
$x$ & $\neg x$ \\
\hline
True & False \\
False & True
\end{tabular}
\end{center}

\textbf{Zweistelliger Operator and:}
\begin{center}
\begin{tabular}{cc|c}
$x$ & $y$ & $x \land y$ \\
\hline
True & True & True \\
True & False & False \\
False & True & False \\
False & False & False
\end{tabular}
\end{center}

\textbf{Zweistelliger Operator or:}
\begin{center}
\begin{tabular}{cc|c}
$x$ & $y$ & $x \lor y$ \\
\hline
True & True & True \\
True & False & True \\
False & True & True \\
False & False & False
\end{tabular}
\end{center}
\end{mainbox}

\begin{subbox}{Priorit¨at \& Assoziativit¨at}
\textbf{Priorit¨at:} not $>$ and $>$ or

\textbf{Assoziativit¨at:} Links nach rechts
\end{subbox}

% ============================================================================
% SECTION 2: INTEGER
% ============================================================================
\section{Integer}

\begin{mainbox}{Definition: Integer}
Ein integer Wert ist eine ganze Zahl:
$$\Z = \{\ldots, -4, -3, -2, -1, 0, 1, 2, 3, 4, \ldots\}$$

\textbf{Operatoren:}
\begin{itemize}
    \item Einstellig: $-$ (Negation)
    \item Zweistellig: $+, -, \ast$ (Addition, Subtraktion, Multiplikation)
\end{itemize}
\end{mainbox}

\begin{mainbox}{Division bei Integer}
\textbf{Ganzzahldivision //}
$$a \mathbin{//} b = \lfloor a/b \rfloor \quad \text{(abgerundet)}$$

\textbf{Rest-Operator \%}
$$a \bmod b = a - b \cdot (a \mathbin{//} b)$$
\end{mainbox}

\begin{mainbox}{Bin¨ardarstellung}
Jede Zahl $n \in \N$ kann als Summe von Zweierpotenzen dargestellt werden:
$$n = \sum_{i=0}^{k} b_i \cdot 2^i \quad \text{mit } b_i \in \{0,1\}$$

\textbf{Algorithmus:}
\begin{align*}
n &= q_0 \cdot 2 + r_0 \quad \text{mit } r_0 \in \{0,1\} \\
q_0 &= q_1 \cdot 2 + r_1 \quad \text{mit } r_1 \in \{0,1\} \\
&\vdots \\
q_{k-1} &= 0 \cdot 2 + r_k \quad \text{mit } r_k \in \{0,1\}
\end{align*}

Bin¨ardarstellung: $(n)_2 = r_k r_{k-1} \ldots r_1 r_0$
\end{mainbox}

% ============================================================================
% SECTION 3: FLOAT
% ============================================================================
\section{Float}

\begin{mainbox}{Definition: Float}
Ein float Wert ist eine Kommazahl mit endlich vielen Nachkommastellen.

\textbf{Zahlenmengen:}
\begin{itemize}
    \item $\Q$: rationale Zahlen (als Bruche darstellbar) ¨
    \item $\R$: reelle Zahlen
\end{itemize}
\end{mainbox}

\begin{mainbox}{Bin¨ardarstellung von Kommazahlen}
Fur Nachkommastellen: ¨
$$x = \sum_{i=1}^{k} b_{-i} \cdot 2^{-i} \quad \text{mit } b_{-i} \in \{0,1\}$$

\textbf{Algorithmus:}
\begin{align*}
x \cdot 2 &= b_1 + r_1 \quad \text{mit } b_1 \in \{0,1\}, r_1 < 1 \\
r_1 \cdot 2 &= b_2 + r_2 \quad \text{mit } b_2 \in \{0,1\}, r_2 < 1 \\
&\vdots
\end{align*}

Bin¨ardarstellung: $(x)_2 = 0.b_1 b_2 b_3 \ldots$

Bei periodischer Wiederholung: $(x)_2 = 0.b_1 b_2 \overline{b_3 b_4}$
\end{mainbox}

\begin{subbox}{IEEE 754 Standard}
\textbf{Fliesskommazahl:} Standard fur float-Darstellung ¨
\begin{itemize}
    \item 32-bit (float in Java)
    \item 64-bit (double in Java)
\end{itemize}

\textbf{Wichtig:} Float-Arithmetik kann unerwartete Ergebnisse haben!
\end{subbox}

% ============================================================================
% SECTION 4: SET, LIST, MATRIX
% ============================================================================
\section{Sammlungen}

\begin{mainbox}{Definition: Menge (Set)}
Eine Menge ist eine \textbf{ungeordnete} Sammlung von Objekten ohne Rucksicht auf Wiederholung. ¨

\textbf{Notation:}
\begin{itemize}
    \item Aufz¨ahlung: $M = \{1, 2, 3, 4\}$
    \item Beschreibung: $M = \{x : x \text{ ist eine Ziffer}\}$
\end{itemize}

\textbf{Elementbeziehung:}
$$x \in M \quad \text{bzw.} \quad x \notin M$$
\end{mainbox}

\begin{mainbox}{Mengenoperationen}
\textbf{Gleichheit:}
$$A = B \iff \forall x: (x \in A \leftrightarrow x \in B)$$

\textbf{Teilmenge:}
$$A \subset B \iff \forall x: (x \in A \rightarrow x \in B)$$

\textbf{Durchschnitt:}
$$A \cap B = \{x : x \in A \land x \in B\}$$

\textbf{Vereinigung:}
$$A \cup B = \{x : x \in A \lor x \in B\}$$

\textbf{Differenz:}
$$A \setminus B = \{x : x \in A \land x \notin B\}$$

\textbf{Komplement:}
$$\overline{A} = \{x : x \notin A\}$$
\end{mainbox}

\begin{mainbox}{Definition: Liste (List)}
Eine Liste ist eine \textbf{geordnete} Sammlung von Objekten mit Rucksicht auf Wiederholung. ¨

\textbf{Notation:}
$$a = (a_1, a_2, \ldots, a_n)$$

\textbf{Indexierung:}
$$a_i \quad \text{oder} \quad a[i] \quad \text{(i-te Komponente)}$$

\textbf{Komponentenweise Operationen:}
$$(a_1, a_2, \ldots, a_n) + (b_1, b_2, \ldots, b_n) = (a_1+b_1, a_2+b_2, \ldots, a_n+b_n)$$

\textbf{Broadcasting:}
$$c \cdot (a_1, a_2, \ldots, a_n) = (c \cdot a_1, c \cdot a_2, \ldots, c \cdot a_n)$$
\end{mainbox}

\begin{mainbox}{Definition: Matrix}
Eine Matrix ist eine \textbf{zweidimensionale}, rechteckig geordnete Sammlung von Objekten.

\textbf{Notation:}
$$A = \begin{pmatrix} a_{11} & a_{12} & \cdots & a_{1n} \\ a_{21} & a_{22} & \cdots & a_{2n} \\ \vdots & \vdots & \ddots & \vdots \\ a_{m1} & a_{m2} & \cdots & a_{mn} \end{pmatrix}$$

\textbf{Dimension:} $m \times n$ (m Zeilen, n Spalten)

\textbf{Indexierung:}
$$a_{ij} \quad \text{oder} \quad A[i,j] \quad \text{oder} \quad A[i][j]$$

\textbf{Komponentenweise Operationen:}
Matrizen gleicher Dimension werden komponentenweise verrechnet.

\textbf{Broadcasting:}
$m \times n$ Matrix mit $1 \times 1$, $m \times 1$ oder $1 \times n$ Matrix.
\end{mainbox}

% ============================================================================
% SECTION 5: FUNCTION
% ============================================================================
\section{Funktionen}

\begin{mainbox}{Definition: Funktion}
Seien $D$ (Definitionsbereich) und $W$ (Wertevorrat) Mengen. Eine Funktion $f$ ordnet jedem Element in $D$ \textbf{genau ein} Element in $W$ zu.

\textbf{Notation:}
$$f : D \longrightarrow W$$
$$x \mapsto y \quad \text{oder} \quad f(x) = y$$
\end{mainbox}

\begin{subbox}{Algorithmus}
Ein \textbf{Algorithmus} ist eine pr¨azise, schrittweise Angabe, die fur jeden Input ¨ $x \in D$ beschreibt, wie der zugeh¨orige Output $y = f(x)$ bestimmt wird.
\end{subbox}

\end{document}